% Generated draft: requires mathematical and source review.
\documentclass[11pt]{article}
\usepackage[margin=25mm]{geometry}
\usepackage{fontspec}
\setmainfont{Latin Modern Roman}
\usepackage{amsmath,amssymb,mathtools,unicode-math}
\setmathfont{Latin Modern Math}
\usepackage{enumitem,needspace,xurl,hyperref}
\setlength{\parindent}{0pt}
\setlength{\parskip}{5pt}
\setlength{\emergencystretch}{3em}
\newcommand{\N}{\mathbb{N}}
\newcommand{\Z}{\mathbb{Z}}
\newcommand{\Q}{\mathbb{Q}}
\newcommand{\R}{\mathbb{R}}
\newcommand{\C}{\mathbb{C}}
\newcommand{\F}{\mathbb{F}}
\newcommand{\A}{\mathbb{A}}
\newcommand{\PP}{\mathbb{P}}
\newcommand{\E}{\mathbb{E}}
\newcommand{\eps}{\varepsilon}
\newcommand{\dd}{\,\mathrm{d}}
\begin{document}
\section*{UnsolvedMath ID: 700103}
\textbf{Source problem code:} catalogue entry
\subsection{Definitions and mathematical statement}
A group embedding \(\iota:G\hookrightarrow H\) is an injective map satisfying \(\iota(xy)=\iota(x)\iota(y)\) for all \(x,y\in G\). Let \(\mathcal O\) denote the class of groups called \emph{orderable} in the source, and let \(\mathcal P\) denote the class of \emph{pro-orderable} groups defined in Kourovka Notebook Problem 1.35.

The question posed by A. I. Kokorin in 1967 is whether
\[
\forall G\in\mathcal O\;\exists H\in\mathcal P\;\exists\iota:G\hookrightarrow H.
\]
Thus the assertion asks for an embedding of each orderable group into some pro-orderable group.

The supplied record does not include the definition from Problem 1.35, so \(\mathcal P\) cannot be explicitly defined from the supplied evidence. It also does not specify the order-invariance convention for \(\mathcal O\): a left-invariant total order satisfies \(x<y\Rightarrow ax<ay\), whereas a two-sided invariant total order additionally satisfies \(x<y\Rightarrow xa<ya\), for all \(a,x,y\in G\). No choice between these conventions is made here. The displayed formulation therefore preserves the source question but remains conditional on recovering its intended definitions.
\subsection{Short English statement}
Does every orderable group admit an injective group homomorphism into a pro-orderable group? The definition of pro-orderability is deferred by the source to Problem 1.35, whose text is absent; the intended convention for orderability also requires confirmation.
\subsection{Sources}
\begin{itemize}
\item Supplied local UnsolvedMath registry record: permanent ID \(700103\), exact source problem code \texttt{catalogue entry}, title \emph{Kourovka Notebook Problem 2.28}. Upstream revision: \texttt{d371ee917ac9647fbb06cd6fc529deec6f3dda82}.
\item \emph{The Kourovka Notebook (Unsolved Problems in Group Theory)}, original Problem 2.28, second issue (1967); proposed by A. I. Kokorin. The supplied record identifies the twenty-first issue edition (2026), arXiv:1401.0300v48: \url{https://arxiv.org/abs/1401.0300v48}. The collection description names E. I. Khukhro and V. D. Mazurov as editors.
\item The statement explicitly refers to Kourovka Notebook Problem 1.35 for the definition of pro-orderable groups; that definition is not included in the supplied record.
\item The registry labels this problem \emph{open} and records an update timestamp of \texttt{2026-10-16T14:00:00Z}. This is supplied registry information, not an independently verified assessment of the literature or mathematical status.
\end{itemize}
\end{document}
